\section{Worked problems: tiles, waves and crossovers}

\subsection{A tile that reaches the ridge}
\textbf{Problem.} Using equation~\eqref{eq:tile-intensity}, find the smallest
square tile $T \times T$ whose intensity in BF16 exceeds an H100's ridge of about
295~FLOP/byte. How much shared memory would one stage of its operands need for a
$K$-slice of 64, and how much room would its accumulators take?

\textbf{Solution.} For a square tile in BF16, $I = 2T^2/(2T \times 2) = T/2$
FLOP/byte, so the tile must be at least $591 \times 591$; say $592 \times 592$. One
stage of operands is $2 \times 592 \times 64 \times 2$~bytes $= 148$~KB --- one
stage fits in the 228~KB of an SM's shared memory, two do not, so the kernel could
not load the next slice while multiplying this one. The accumulators are
$592^2 \times 4$~bytes $= 1.34$~MiB in 32-bit floats: six times the SM's entire
shared memory, and far more than its registers can hold. No single SM's tile
reaches the ridge from HBM. Real kernels get there by reusing tiles \emph{across}
SMs: neighbouring blocks read the same rows of $\mathbf{A}$ or columns of
$\mathbf{B}$ in turn, so most of their loads hit the 50~MB L2 cache rather than
HBM, and on Hopper a cluster of up to 16 blocks can read one another's shared
memory directly \cite{nvidia2026hoppertuning}. The ridge is a property of the
whole chip's reuse, not of one tile.

\subsection{Recovering a lost wave}
\textbf{Problem.} A product divides into 150 output tiles on a GPU with 132 SMs,
one tile per SM at a time. Compare its time with that of a 132-tile product, and
evaluate two remedies: splitting every tile's $K$ loop into equal parts
(split-K), and stream-K, which gives each SM an equal share of all the tiles'
$K$ iterations.

\textbf{Solution.} Without a remedy the product takes two tile-times --- 132 tiles
in the first wave, 18 in the second --- and the SMs do useful work in
$150/264 = 57\%$ of the slots. Splitting each tile's $K$ loop in two gives 300
half-tiles: three waves of half-tiles, 1.5 tile-times, 75\% of the slots used,
plus a reduction that adds 150 pairs of partial tiles. The best uniform split,
into sevenths, gives 1,050 units in eight waves of a seventh each: 1.14
tile-times, but with seven partial tiles to add for every output tile. Stream-K
gives each SM $150/132 = 1.14$ tiles' worth of $K$ iterations directly: 1.14
tile-times, with every SM busy until the end, and only the tiles whose
iterations are divided between two SMs need a fix-up of partial sums (derived).
Both remedies turn two tile-times into about 1.14; stream-K pays for it with
fewer reductions, which is why its authors could use one tile configuration for
thousands of shapes \cite{osama2023streamk}.

\subsection{Predicting the crossover}
\textbf{Problem.} In the lab, the streaming decode kernel takes 1.2~ms at small
$N$ (the 67~MB weight matrix read at about 57~GB/s) and computes at about
110~GFLOP/s once the arithmetic dominates; the tiled GEMM takes 3.28~ms for any
$N$ from 1 to 16. Predict the $N$ at which the tiled kernel becomes faster, and
compare with the measurement.

\textbf{Solution.} The streaming kernel's time is the longer of 1.2~ms and
$2 \times 4{,}096^2 \times N / 110$~GFLOP/s $= 0.305N$~ms. It reaches the tiled
kernel's 3.28~ms at $N = 3.28 / 0.305 \approx 11$ (derived). Measured, the
streaming kernel took 2.2~ms at $N = 8$ and 6.0--7.2~ms at $N = 16$ against the
tiled kernel's 3.3~ms at both, so the crossover lies between 8 and 16, as
predicted. The prediction is slightly optimistic at 16 --- it gives 4.9~ms ---
because the streaming kernel's arithmetic rate falls once sixteen vectors no
longer fit beside a weight row in the first-level cache. An engine's switching
threshold is this crossover, measured for its own kernels on its own hardware;
llama.cpp's Metal backend switches at eight rows (\Chref{roofline}), and a
different machine or kernel would move it.
